\documentclass{ifcolog}

\usepackage{makeidx}
\usepackage{amsthm}
\usepackage{amssymb}
\usepackage{algorithm}
\usepackage{algpseudocode}

\newcommand{\bc}{ERG}
\newcommand{\stc}{stablecoin}
\newcommand{\Pdt}{P_{\mathrm{D}}^{*}}
\newcommand{\Plp}{P_{\mathrm{D}}^{\mathrm{LP}}}
\newcommand{\Sd}{S_{\mathrm{D}}}
\newcommand{\rr}{\bar{r}}

\newcommand{\sbank}{s_{\textrm{BANK}}}
\newcommand{\slp}{s_{\mathrm{LP}}}
\newcommand{\soracle}{s_{\textrm{O}}}
\newcommand{\sfree}{s_{\textrm{FREE}}}
\newcommand{\sarb}{s_{\textrm{ARB}}}
\newcommand{\sib}{s_{\textrm{IB}}}
\newcommand{\sburn}{s_{\textrm{BURN}}}

\newcommand{\tautx}{\tau_{\textrm{TX}}}

\newcommand{\epsilonfree}{\epsilon_{\textrm{FREE}}}
\newcommand{\epsilondexyzero}{\epsilon_{\textrm{ED}}^{0}}
\newcommand{\epsilondexyone}{\epsilon_{\textrm{ED}}^{1}}
\newcommand{\epsilonreleasedexy}{\epsilon_{\textrm{RD}}}

\newcommand{\dx}{Dexy}
\newcommand{\dxg}{DexyGold}
\newcommand{\dxusd}{USE}

\newtheorem{definition}{Definition}
\newtheorem{theorem}{Theorem}

\title{Dexy:\protect \\ A Simple Stablecoin Design Based on an Algorithmic Central Bank}

\titlerunning{Dexy}

%\addauthor{IFCoLog}{International Federation of Computational Logic}
\addauthor[kushti@protonmail.ch]
  {Alexander Chepurnoy}
  {Ergo Platform}
\addauthor[]
  {Amitabh Saxena}
  {Ergo Platform}
\addauthor[ldgaetano@protonmail.com]
  {Luca D'Angelo}
  {The Stable Order}

\authorrunning{A. Chepurnoy and A. Saxena and L. D'Angelo}

\titlethanks{
We would would like to thank Ile for his inspiring forum posts, 
Bruno for discussions and the extract to future idea, and  4EYES Labs for 
the audit of the ErgoScript contracts.
}

\begin{document}
\maketitle

\begin{abstract}
We consider a new algorithmic stablecoin called Dexy, consisting of two assets: the base cryptocurrency, called the basecoin, and the stablecoin.
The protocol consits of three components: a central bank, a secondary market in the form of a customized CP-AMM liquidity pool, and a trusted oracle. 
The price of the stablecoin in the liquidity pool is stabilized, relative to the oracle price, via central bank interventions. 
\end{abstract}

\section{Introduction}
\label{sec:introduction}
Due to the inherent volatility of cryptocurrency prices, algorithmic stablecoins naturally arose as their extension. The volatility is resolved 
by pegging the stablecoin to an asset whose price is considered to be stable over long periods of time, e.g. gold. 

Having an asset with a stable value can be useful in many scenarios:
\begin{itemize}
\item Securing fundraising: a project can be sure that funds collected during fundraising will have stable value in the mid-to-long term.
\item Doing business with predictable results: a shop can be sure that funds collected from sales will have roughly the same value when the shop is ordering goods from warehouses, otherwise, the shop may go bankrupt if its margin is to low to cover exchange rate fluctuations. 
\item Shorting: when cryptocurrency prices are high, it is desirable for investors to rebalance their portfolio by increasing exposure to fiat currencies or real-world commodities. However, as fiat currencies and centralized exchanges impose significant risks, it would be better to buy fiat and commodity substitues in the form of stablecoins on decentralized exchanges.
\item Lending and other decentralized finance applications: stability of collateral value is critical for many applications built on top of these services.
\end{itemize}

Centralized stablecoins, such as USDT and USDC, use a trusted party to ensure that the peg is maintained. For algorithmic stablecoins, the 
peg can be maintained in different ways depending on the implementation, e.g. rebasement of the total supply \cite{Ampleforth}, 
or via imitating a trusted party that holds reserves and then does market interventions for rebalancing the peg. 
Imititating the trusted party is usually done by allowing anyone on the blockchain to create over-collateralized financial instruments, 
such as collateralized debt positions \cite{DAI}, zero-coupon bonds (as was done in the failed Yield protocol), reserve assets \cite{DjedEprint,DjedICBC2023,Gluon}, 
or by issuing stabilizing financial instruments in case of a depeg \cite{Neutrino}.

In this work, we present \dx{}, a stablecoin protocol where an algorithmic central bank performs interventions when there is a depeg.
The bank stabilizes the stablecoin value in the reference market, a liquidity pool, injecting basecoins from its reserves when the stablecoin price is below the peg. 
In extreme cases, when bank reserves are depleted and the stablecoin price is still below the peg, the liquidity pool price is restored by extracting stablecoins, 
which will be injected back again at a later time when the stablecoin price is above the peg by a certain threshold. 
In some aspects, \dx{} could resemble Fei or Gyroscope \cite{Gyroscope}. 

The rest of the paper is organized as follows.  In Section \ref{sec:design}, the \dx{} design in general is explained, where important trust assumptions are stated and
notation used in other sections is introduced. Section \ref{sec:worst-case} defines the worst-case scenario for the bank reserves, which implies for the stablecoin stability as well. 
In Section \ref{sec:rules}, the bank and liquidity pool rules are defined, where their stability is discussed in Section \ref{sec:stability}. 
Section \ref{sec:implementations} describes \dx{} implementations: the gold-pegged DexyGold stablecoin and the USD-pegged USE stablecoin.

\section{\dx Design}
\label{sec:design}

The \dx protocol functions via an algorithmic central bank that holds reserves in basecoins and mints stablecoins. 
A corresponding liquidity pool exists where stablecoins and basecoins can be traded. 
The peg price used by the bank is provided by an oracle, which is assumed to be trusted and reliable. 
Users can interact with both the bank and the liqudity pool, but certain rules and restrictions are applied to control the price stability.
Central bank interventions are determined by the ratio of the liquidity pool price to the oracle price, $r$, which we also call the reference reserve ratio. 

%TODO: Insert fig from workshop presentation.

\subsection{The Liquidity Pool}
\label{subsec:pool}

The reference market is implemented as an automated market maker (e.g. \cite{xu2021sok}). 
For simplicity, a modified constant-product automated market maker (CP-AMM) is used, similar to ErgoDex \cite{ErgoDex} and UniSwap \cite{Uniswap}. 
This means that the liqudity pools holds a pair of basecoins and stablecoins such that the product of their amounts remains invariant before and after a trade. 
Users who function as liquidity providers can provide a pair of basecoins and stablecoins to obtain a corresponding amount of LP tokens.
These LP tokens can later be deposited to obtain the initial liquidity pair deposit plus accumulated fees from trading. 
A fee is charged to liqudity providers to prevent excessive liquidity withdrawals.

%TODO: Insert fig from workshop presentation.

\subsection{The Algorithmic Central Bank}
\label{subscec:bank}

By depositing basecoins into the bank, users can mint stablecoins. 
The price provided by the bank will always equal the oracle price. 
A fee is charged so the bank reserves can increase faster.
Under certain conditions, the bank can intervene into the liquidity pool to control its price relative to the oracle price.

%TODO: Insert fig from workshop presentation.

%TODO: List protocol actions and who triggers them.

\subsection{Trust Assumptions}
\label{subsec:kya}

In general, for any DeFi protocol involving user funds, it is important for trust assumptions to be explicitly stated so that users can make informed decisions to the best of their ability.
The \dx protocol has two:

\begin{itemize}
  \item{Oracle:} The protocol assumes the target price oracle is reliable. This trust assumption is unavoidable; however, it can be relaxed by using a federation of oracles that produces an average price 
  to be delivered to the central bank, instead of relying on a single entity. In any case, it is important that oracles are rewarded from protocol usage, 
  otherwise, they will be incentivized to adversarially attack the protocol, see Section \ref{subsec:oracle} for details.

  \item{No governance:} Unlike most stablecoin protocols, we do not consider governance, as it can be used to adversarially attack the protocol. 
  We suppose that real-world instantiations that do have it should clearly state their governance-related trust assumptions.
\end{itemize}

\subsection{Actions and Dynamics}
\label{subsec:actions}

The following table shows the different actions in the protocol, the $r$ range when the action happens, and the restored $r'$ range after the action happens.

\begin{table}[h]
\begin{tabular}{|l|l|l|}
\hline
\multicolumn{1}{|c|}{\textbf{Action}} & \multicolumn{1}{c|}{\textbf{Range}}                                   & \multicolumn{1}{c|}{\textbf{Restore Range}}                                                         \\ \hline
FreeMint                              & $1- \epsilon_{\textrm{FREE}} \le r \le 1+ \epsilon_{\textrm{FREE}}$   & \multicolumn{1}{c|}{-}                                                                              \\ \hline
ArbitrageMint                         & $r > 1 + \epsilon_{\textrm{FREE}}$                                    & \multicolumn{1}{c|}{-}                                                                              \\ \hline
InjectBasecoins                       & $r < 1 - \epsilon_{\textrm{FREE}}$                                    & $1 - \epsilon_{\textrm{FREE}} \le r' \le 1-\epsilon_{\textrm{IB}}$                                  \\ \hline
ExtractDexy                           & $r \le 1 - \epsilon_{\textrm{ED}}^{0} < 1 - \epsilon_{\textrm{FREE}}$ & $1 - \epsilon_{\textrm{ED}}^{0} + \epsilon_{\textrm{ED}}^{1} \le r' \le 1-\epsilon_{\textrm{FREE}}$ \\ \hline
ReleaseDexy                           & $1 + \epsilon_{\textrm{RD}} < r <  1 + \epsilon_{\textrm{FREE}}$      & $r' \ge 1 + \epsilon_{\textrm{RD}}$                                                                  \\ \hline
\end{tabular}
\end{table}

We now go over at a high-level what happens during each action. For more technical details, refer to section $\ref{sec:protocol}$.

\subsubsection{FreeMint}
  \begin{itemize}
    \item  
  \end{itemize}
\subsubsection{ArbitrageMint}
\subsubsection{InjectBasecoins}
\subsubsection{ExtractDexy}
\subsubsection{ReleaseDexy}

\section{Protocol Specification}
\label{sec:protocol}
We now describe the \dx{} protocol specification, which includes immutable paramters and state variables.
% \begin{itemize}
%   \item{} $T$ - period before intervention starts. After one intervention, the bank can start another one only after $T$ time units has passed. 
%   \item{} $p$ - basecoin price reported by the oracle at the moment (e.g. 20 USD per ERG, note, however, that the price provided by the oracle will usually be in units of $\frac{\mathrm{basecoin}}{\mathrm{stablecoin}}$, representing the price of the stablecoin in units of basecoins).
%   \item{} $s$ - price when the bank should intervene in case of a sudden basecoin price crash. For example, we can assume that $s = \frac{p}{4}$ (so if p is 20 USD per ERG, then $s$ is 5 USD per ERG, means the bank needs to have enough reserves to save the markets when the price is suddenly crashing from 20 to 5 USD per ERG).
%   \item{} $R$ - ratio between $p$ and $s$, $R = \frac{p}{s}$.
%   \item{} $r$ - ratio between $p$, and price in the pool, which is $\frac{u}{e}$, thus $r = \frac{p}{\frac{u}{e}} = \frac{p*e}{u}$.
%   \item{} $e$ - amount of basecoin in the liquidity pool. 
%   \item{} $u$ - amount of \stc{} in the liquidity pool.
%   \item{} $O$ - amount of \stc{} in circulation that is outside the liquidity pool. The amount $O$ is not known for the \dx{} protocol, but the bank can easily store how many \stc{} tokens were minted and then get $O$ by deducting $u$ from it.
%   \item{} $E$ - amount of basecoins in the bank. 
% \end{itemize}

% Let A be the set of protocol actions such that:

% $$A = \left\{ \textrm{ARB}, \textrm{FREE}, \textrm{IB}, \textrm{ED}, \textrm{RD} \right\}$$

% \begin{itemize}
%   \item ARB is the arbitrage-mint action.
%   \item FREE is the free-mint action.
%   \item IB is the inject-basecoins action.
%   \item ED is the extract-dexy action.
%   \item RD is the release-dexy action.
% \end{itemize}

The central bank parameters is a tuple $C_{\textrm{BANK}}$ of:
$$\left( \phi_{\textrm{BANK}}^{\textrm{FREE}}, \phi_{\textrm{BANK}}^{\textrm{ARB}}, \epsilon_{\textrm{FREE}}, \epsilon_{\textrm{IB}}, \epsilon_{\textrm{ED}}^{0}, \epsilon_{\textrm{ED}}^{1}, \epsilon_{\textrm{RD}}, T_{\textrm{FREE}}, T_{\textrm{ARB}}, T_{\textrm{IB}},  \tau_{\textrm{IB}}, \tau_{\textrm{ED}}, \tau_{\textrm{RD}}, \tau_{\textrm{TX}}, \alpha_{\textrm{FREE}}, \alpha_{\textrm{IB}}  \right)$$

\begin{itemize}
  \item $\phi_{\textrm{BANK}}^{\textrm{FREE}} \in (0, 1)$ is the bank fee for the free-mint action.
  \item $\phi_{\textrm{BANK}}^{\textrm{ARB}} \in (0,1)$ is the bank fee for the arbitrage-mint action.
  \item $\epsilon_{\textrm{FREE}} \in (0, 1)$ defines the upper and lower bounds of the free-mint range, representing the percentage that the LP price deviates from the oracle price.
  \item $0 < \epsilon_{\textrm{IB}} < \epsilon_{\textrm{FREE}}$ defines the percentage where the inject-basecoins action will restore the price.
  \item $\epsilon_{\textrm{FREE}} < \epsilon_{\textrm{ED}}^{0} < 1$ defines the percentage where the extract-dexy action will happen.
  \item $0 < \epsilon_{\textrm{ED}}^{1} < \epsilon_{\textrm{ED}}^{0}$ defines the percentage where the extract-dexy action will restore the price.
  \item $\epsilon_{\textrm{RD}} < \epsilon_{\textrm{FREE}}$ defines the percentage where the release-dexy action will happen.
  \item $T_{\textrm{FREE}} \in \mathbb{R}_{>0}$ is the free-mint period.
  \item $T_{\textrm{ARB}} \in \mathbb{R}_{>0}$ is the arbitrage-mint period.
  \item $T_{\textrm{IB}} \in \mathbb{R}_{>0}$ is the inject-basecoins period.
  \item $\tau_{\textrm{IB}} \in \mathbb{R}_{>0}$ is the inject-basecoins bank delay.
  \item $\tau_{\textrm{ED}} >> \tau_{\textrm{IB}}$ is the extract-dexy bank delay, which must be much larger than the inject-basecoins bank delay, since extraction is an action of last-resort.
  \item $\tau_{\textrm{RD}} \in \mathbb{R}_{>0}$ is the release-dexy bank delay.
  \item $\tau_{\textrm{TX}} \in \mathbb{R}_{>0}$ is the allowed transaction propogation delay.
  \item $\alpha_{\textrm{FREE}} \in (0,1)$ is the percentage of dexy liquidity, in the liquidity pool, that determines the max amount of dexy that can be minted from the bank during the free-mint period.
  \item $\alpha_{\textrm{IB}} \in (0,1)$ is the percentage of the bank's basecoin reserve that determines the max amount of basecoins to inject into the liquidity pool during an inject-basecoin action.
\end{itemize}

The central bank state is determined by a tuple of two values:
$$s_{\textrm{BANK}} = \left(R, S_{D} \right)$$

\begin{itemize}
  \item $R \in \mathbb{R}_{>0}$ is the bank's basecoins reserves.
  \item $S_D \in \mathbb{R}_{>0}$ is the total supply of minted dexy stablecoins.
\end{itemize}

The liquidity pool has only one parameter:
$$C_{LP} = \left(\phi_{LP}\right)$$

\begin{itemize}
  \item $\phi_{LP} \in (0,1)$ is the liquidity provider withdrawal fee, which should be high enough to disincentivize frequent withdrawals. 
\end{itemize}

The liquidity pool state is determined by a tuple of three values:
$$s_{\textrm{LP}} = \left(B, D, S_{LP} \right)$$

\begin{itemize}
  \item $B \in \mathbb{R}_{>0}$ is the amount of basecoins in the liquidity pool.
  \item $D \in \mathbb{R}_{>0}$ is the amount of dexy stablecoins in the liquidity pool.
  \item $S_{LP} \in \mathbb{R}_{>0}$ is the supply of LP tokens provided to liquidity providers.
\end{itemize}

The oracle parameters are assumed to only be the following:
$$C_{O} = \left(\phi_{O}\right)$$

\begin{itemize}
  \item $\phi_{O} \in (0,1)$ is the oracle fee.
\end{itemize}

The oracle state is determined by a tuple of two values:
$$s_{\textrm{O}} = \left(\Pdt, F \right)$$

\begin{itemize}
  \item $\Pdt \in \mathbb{R}_{>0}$ is the target stablecoin price, in units of $\frac{\textrm{BASECOIN}}{\textrm{PEG ASSET}}$.
  \item $F \in \mathbb{R}_{\ge 0}$ is the accumulated fees of the oracle.
\end{itemize}

For the sake of clarity, we now define separate state variables for the different bank actions, 
but these can also be considered as part of the bank state.

The FreeMint action has the following state:
$$s_{\textrm{FREE}} = \left( t, d \right)$$

\begin{itemize}
  \item $t \in \mathbb{R}_{>0}$ is the time when the free-mint period ends.
  \item $d \in \mathbb{R}_{\ge0}$ is the remaining amount of dexy stablecoins that can be minted during the free-mint period.
\end{itemize}

The ArbMint action has the following state:
$$s_{\textrm{ARB}} = \left(t, d, \tau\right)$$

\begin{itemize}
  \item $t \in \mathbb{R}_{>0}$ is the time when the arbitrage-mint period ends.
  \item $d \in \mathbb{R}_{\ge 0}$ is the remaining amount of dexy stablecoins that can be minted during the arbitrage-mint period.
  \item $\tau \in \mathbb{R}_{>0}\cup\{\infty\}$ is the time when $r > 1 + \epsilon_{\textrm{FREE}}$. 
\end{itemize}

The InjectBasecoin action has the following state:
$$s_{\textrm{IB}} = \left(t, \tau\right)$$

\begin{itemize}
  \item $t \in \mathbb{R}_{>0}$ is the time when the previous inject-basecoin action occurred.
  \item $\tau \in \mathbb{R}_{>0}\cup\{\infty\}$ is the time when $r < 1 + \epsilon_{\textrm{FREE}}$. 
\end{itemize}

The ExtractDexy and ReleaseDexy actions share the following state:
$$s_{\textrm{BURN}} = \left(t_{\textrm{ED}}, t_{\textrm{RD}}, d, \tau_{\textrm{ED}}, \tau_{\textrm{ED}} \right)$$

\begin{itemize}
  \item $t_{\textrm{ED}} \in \mathbb{R}_{>0}$ is the time when the previous extract-dexy action occurred.
  \item $t_{\textrm{RD}} \in \mathbb{R}_{>0}$ is the time when the previous release-dexy action occurred.
  \item $d \in \mathbb{R}_{\ge 0}$ is the remaining dexy amount to be released.
  \item $\tau_{\textrm{ED}} \in \mathbb{R}_{>0}\cup\{\infty\}$ is the time when $r \le 1 - \epsilon_{\textrm{ED}}^{0}$.
  \item $\tau_{\textrm{RD}} \in \mathbb{R}_{>0}\cup\{\infty\}$ is the time when $r > 1 + \epsilon_{\textrm{RD}}$.  
\end{itemize}

\subsection{Auxiliary Definitions}

\begin{definition}[Bank Liabilities]
  The bank liabilities is defined as the product of 
  the oracle price and the minted supply of dexy stablecoins:
  \begin{equation}
    L(\Pdt, \sbank) = \Pdt \Sd
  \end{equation}
\end{definition}

\begin{definition}[Bank Reserve Ratio]
  The bank reserve ratio is defined as the ratio of its reserves to its liabilities:
  \begin{equation}
    r(\Pdt, \sbank) = \frac{R}{L(\Pdt, \sbank)} = \frac{R}{\Pdt \Sd}
  \end{equation}
  Assuming the oracle price is constant, and the reserves are greater than the liabilities,
  the bank reserve ratio will always increase after stablecoins are minted. 
  During periods of sudden basecoin price crashes, the reserve ratio will decrease. 
  Since users cannot sell dexy stablecoins to the bank, the bank is liable to no one, 
  thus the bank reserves can be less than its liabilities, which is expected during bank interventions.
  So, a measure of the health of the $\dx{}$ protocol is having an increasing reserve ratio, 
  before and after bank interventions.
 
\end{definition}

\begin{definition}[Reference Price]
  The liquidity pool price, or the reference market price, is defined as follows:
  \begin{equation}
    \Plp(\slp) = \frac{B}{D}
  \end{equation}
  The reference price, $\Plp$, has units of $\frac{\textrm{BASECOIN}}{\textrm{PEG ASSET}}$.
\end{definition}

\begin{definition}[Reference Reserve Ratio]
    The ratio of the liquidity pool price to the oracle price is called the reference reserve ratio:
\begin{equation}
  \label{eq:r}
  \bar{r}(\Pdt, \slp) = \frac{\Plp(\slp)}{\Pdt} = \frac{B}{\Pdt D}
\end{equation}
  The aim of the $\dx{}$ protocol is to maintain $\rr$ within a sufficiently small neighbourhood of 1.
  $\rr$ is called the reference reserve ratio since we can interpret the formula as the reserve ratio of the liquidity pool.
\end{definition}

\subsection{State Update Procedures}

\subsubsection{UpdateArbitrageMint}

\begin{algorithm}
\caption{UpdateArbitrageMint}\label{alg:update-arbmint}
\begin{algorithmic}[1]
\State $(\tilde{t}, d, \tau) \gets \sarb$
\State $\rr \gets \rr(\Pdt, \slp) $
\State $notUpdated \gets (\tau == \infty)$
\State $validUpdate \gets (t' \ge t - \tautx) \land (t' \le t)$ \Comment{Tx submission time is in a valid interval.}
\State $notReset \gets (\tau < \infty)$
\State $validReset \gets (t' == \infty)$
\State $isUpdate \gets (\rr > 1  + \epsilonfree) \land notUpdated \land validUpdate$ 
\State $isReset \gets (\rr \le 1 + \epsilonfree) \land notReset \land validReset$
\If{$isUpdate \lor isReset$}
  \State $\sarb \gets (\tilde{t}, d, t')$
\Else
  \State $\sarb \gets \sarb$
\EndIf
\end{algorithmic}
\end{algorithm}

% \begin{equation}
% \operatorname{UpdateArbitrageMint}(\Pdt,\slp,\sarb, t, t'):
% \begin{cases}
%   (\tilde{t}, d, \tau) \leftarrow \sarb \\
%   \rr \leftarrow \rr(\Pdt, \slp) \\
%   \textrm{isUpdate} \leftarrow (\rr > 1 + \epsilonfree) \\
%   \textrm{isReset} \leftarrow (\rr \le 1 + \epsilonfree) \\
%   if 
%   % \textrm{notUpdated} \leftarrow (\tau == \infty) \\
%   % \textrm{isUpdate} \leftarrow (t' \geq t - \tautx) \land (t' \le t) \\
%   % \textrm{notReset} \leftarrow (\tau < \infty) \\
%   % \textrm{isReset} \leftarrow (t' == \infty) \\
%   % \textrm{update} \leftarrow (\rr > 1 + \epsilonfree) \land notUpdated \land isUpdate  \\
%   \text{if } t \ge T_{\mathrm{dep}}(u) + \Delta_{\mathrm{dep}} \text{ and } L_\ell(u) > 0 \\
%     \quad L_f(u) \leftarrow L_f(u) + L_\ell(u) \\
%     \quad L_\ell(u) \leftarrow 0.
% \end{cases}
% \label{eq:update-arbmint}
% \end{equation}

\subsubsection{UpdateInjectBasecoins}

\begin{algorithm}
\caption{UpdateInjectBasecoins}\label{alg:update-inject-basecoins}
\begin{algorithmic}[1]
\State $(\tilde{t}, d, \tau) \gets \sarb$
\State $\rr \gets \rr(\Pdt, \slp) $
\State $notUpdated \gets (\tau == \infty)$
\State $validUpdate \gets (t' \ge t - \tautx) \land (t' \le t)$ \Comment{Tx submission time is in a valid interval.}
\State $notReset \gets (\tau < \infty)$
\State $validReset \gets (t' == \infty)$
\State $isUpdate \gets (\rr < 1  - \epsilonfree) \land notUpdated \land validUpdate$ 
\State $isReset \gets (\rr \ge 1 - \epsilonfree) \land notReset \land validReset$
\If{$isUpdate \lor isReset$}
  \State $\sarb \gets (\tilde{t}, d, t')$
\Else
  \State $\sarb \gets \sarb$
\EndIf
\end{algorithmic}
\end{algorithm}

\subsubsection{UpdateExtractDexy}

\begin{algorithm}
\caption{UpdateExtractDexy}\label{alg:update-extract-dexy}
\begin{algorithmic}[1]
\State $(\tilde{t}, d, \tau) \gets \sarb$
\State $\rr \gets \rr(\Pdt, \slp) $
\State $notUpdated \gets (\tau == \infty)$
\State $validUpdate \gets (t' \ge t - \tautx) \land (t' \le t)$ \Comment{Tx submission time is in a valid interval.}
\State $notReset \gets (\tau < \infty)$
\State $validReset \gets (t' == \infty)$
\State $isUpdate \gets (\rr \le 1 - \epsilondexyzero) \land notUpdated \land validUpdate$ 
\State $isReset \gets (\rr > 1 - \epsilondexyzero) \land notReset \land validReset$
\If{$isUpdate \lor isReset$}
  \State $\sarb \gets (\tilde{t}, d, t')$
\Else
  \State $\sarb \gets \sarb$
\EndIf
\end{algorithmic}
\end{algorithm}

\subsubsection{UpdateReleaseDexy}

\begin{algorithm}
\caption{UpdateReleaseDexy}\label{alg:update-release-dexy}
\begin{algorithmic}[1]
\State $(\tilde{t}, d, \tau) \gets \sarb$
\State $\rr \gets \rr(\Pdt, \slp) $
\State $notUpdated \gets (\tau == \infty)$
\State $validUpdate \gets (t' \ge t - \tautx) \land (t' \le t)$ \Comment{Tx submission time is in a valid interval.}
\State $notReset \gets (\tau < \infty)$
\State $validReset \gets (t' == \infty)$
\State $isUpdate \gets (\rr > 1  + \epsilonreleasedexy) \land notUpdated \land validUpdate$ 
\State $isReset \gets (\rr \le 1 + \epsilonreleasedexy) \land notReset \land validReset$
\If{$isUpdate \lor isReset$}
  \State $\sarb \gets (\tilde{t}, d, t')$
\Else
  \State $\sarb \gets \sarb$
\EndIf
\end{algorithmic}
\end{algorithm}

\subsection{Protocol Operations}

\subsubsection{FreeMint}

\subsubsection{ArbitrageMint}

\subsubsection{InjectBasecoins}

\subsubsection{ExtractDexy}

\subsubsection{ReleaseDexy}

\section{Worst Case Scenario and Bank Reserves}
\label{sec:worst-case}

Possible failure scenarios for such a design, and how to put restrictions and design rules for the system to ensure stability of the \stc{} token price, will now be considered. 

The bank intervenes when there is a sudden drop in the basecoin price provided by the oracle and enough time has passed for the markets to have stabilize themselves, with no interventions, but have failed to do so. In our case, the bank is doing interventions based on the stablecoin price in the liquidity pool in comparison with the oracle price. So, the bank's intervention is about injecting its basecoin reserves into the pool, thereby decreasing the basecoin price to match the oracle price.  

Consider an example where the basecoins are \bc{}, and assume that the oracle price crashes from $p$ to $s$ sharply and stands there, and before the crash there were $e$ of \bc{} and $u$ of \stc{} in the liquidity pool, with pool's price being $p$. The worst case then is when no liqudity is put into the pool during the period $T$. With large enough $T$ and large enough $R$, this assumption is not very realistic: some traders will buy \bc{}s with their \stc{}s anyway, price is usually failing with swings where traders could mint \stc{}s thus increasing \bc{} bank reserves, etc. However, it would be reasonable to consider worst-case scenario, then in the real world \dx{} will be even more durable than in theory. 

In this case, the bank must intervene after $T$ units of time, as the price differs significantly, and restore the price in the pool, so set it to $s$. We denote amounts of basecoins and \stc{}s in the pool after the intervention as $e'$ and $u'$, respectively. Then:

\begin{itemize}
  \item{} $e * u = e' * u'$.
  \item{} The bank must inject $E_1$ basecoins into the pool, thus after the trade there will be $e' = e + E_1$ basecoins in the pool.
  \item{} The new basecoin price in the pool will be $\frac{u'}{e'} = s$, thus $u' = s * (e + E_1)$.
  \item{} So, isolating for the amount that must be injected, we obtain $E_1 = \sqrt{\frac{e * u}{s}} - e$ basecoins that must be added into the pool from the bank.
\end{itemize}

So by injecting $E_1$ basecoins, the bank recovers the price. However, this is not enough, as now there are $O$ \stc{} units which can be injected into the pool from outside. 
Again, it is not realistic to assume that all the $O$ \stc{} would be injected, as some of them are simply lost, some would be kept to buy cheap basecoins at the bottom~(\stc{}s are often used as a bet for a basecoin price decline), etc. However, we need to assume the worst-case scenario again. We also unrealistically assume that all the $O$ tokens are being sold in very small batches that do not significantly affecting the pool price, where after each batch, a seller of a new batch is waiting for a bank intervention to happen (so for $T$ units of time), and sells only after the intervention. In this case, all the $O$ tokens are being sold at price close to $s$, so the bank should have about $E_2 = \frac{O}{s}$ basecoins in reserve to buy the tokens back.

Summing $E_1$ and $E_2$, we obtain the basecoin reserves the bank should have to be ready for the worst-case scenario: $E_w = E_1 + E_2 = \sqrt{\frac{e * u}{s}} - e + \frac{O}{s}$.

This scenario is the worst-case for the bank~(and only the bank) because it is buying all of the external $O$ $\stc{}$ at the worst possible price $s$, and also sets the new price by burning its own reserves only.  

\section{Bank and Pool Rules}
\label{sec:rules}

Based on needed reserves for worst-case scenario estimation, we can consider minting rules accordingly. Similarly to SigUSD \cite{SigUSD} (a Djed instantiation on the Ergo blockchain), we can, for example, target for security in case of a
4x price drop, so to consider $R = \frac{p}{s} = 4$. 
We can also allow minting of \stc{}s~while there are enough, so not less than $E_w$, basecoins in reserves. However, in this case, the \stc{} would be non-mintable most of the time, and only during moments of basecoin price going up significantly will it be possible to mint \dx{}. As worst-case scenario is based on unrealistic assumptions, unlikely a realistic protocol can be built on top of it.  

Thus, we leave worst-case scenario for UI, so dapps implemented to interact with the \dx{} protocol may show, for example, collateralization for the worst-case scenario. Having on-chain data analysis, more precise estimations of reserve quality can be made~(by considering \stc{}s locked in DeFi protocols, stablecoins likely forgotten in wallets, etc).

Instead, we always allow for cautious minting. That is, we allow minting when oracle price is above pool's price~(so the bank provides liquidity for arbitrage when the \stc{} price is above the peg), or we allow to mint a little bit (per some time period) when liquidity pool is in good shape. In details, we have two following minting operations, with minting price being the oracle's price $p$:  

\begin{itemize}
  \item {Arbitrage mint:} if price reported by the oracle is higher than in the pool, i.e. $p > \frac{u}{e}$, we allow to print enough \stc{} tokens to bring the price to $p$. That is, the bank allows to mint up to $\delta_u = \sqrt{p*e*u}-u$ \stc{} by accepting up to $\delta_e = \frac{\delta_u}{p}$ basecoins into reserves. 
  
    \emph{To instantiate the rule, we can allow for arbitrage minting if the price $p$ is more than $\frac{u}{e}$ by at least $1\%$ for time period $T_{arb}$ (e.g. 1 hour), with the bank charging $0.5\%$ fee for the operation. After arbitrage mint, it is not possible to do another one within 30 minutes, in order to prevent aggressive liquidity minting via chained transactions, etc.}

  \item {Free mint:} we allow to mint up to $\frac{u}{100}$ \stc{}s within time period $T_{free}$. 

    \emph{To instantiate the rule, we propose to allow for free mint if $0.98 < r < 1.02$. Minting fee in this case is also $0.5\%$. We propose to set $T_{free}$ to $12$ hours, then bank reserves can grow by $2\%$ of LP volume per day when the peg is okay.}
\end{itemize}

In addition to minting actions, which increase the bank reserves, we define the following action which decreases the bank reserves, but functions as the first stabilization mechanism: 

\begin{itemize}
   \item{Intervention:} if the price reported by the oracle is lower than the pool price by a significant enough margin, i.e. $r = \frac{p*e}{u}$ is less than or equal to some constant that is set as a protocol parameter, then the bank will provide enough basecoins to restore the pool price to the oracle price.
   
   \emph{To instantiate the rule, we propose to allow the bank to intervene if $r \le 0.98$ for time period $T_{int} = 1 {\ day}$.
       During the intervention, the period is reset, so another intervention will only occur after $T_{int} = 1 {\ day}$ at least. The bank will provide enough basecoins to get the pool price
       price to 99.5\% of the oracle price max, but making sure to not spend more than 1\% of its reserves for that.}
\end{itemize}

We also state following rules for the liquidity pool (which, otherwise, acts as an ordinary CP-AMM liquidity pool): 

\begin{itemize}
   \item{Stopping withdrawals:} if $r$ is below some threshold, withdrawals from liquidity providers are stopped (i.e. liquidity providers cannot redeem their LP tokens for basecoins/stablecoins), so only swaps between basecoins and stablecoins are possible.  
   
   \emph{To instantiate the rule, we propose to stop withdrawals immediately if $r \le 0.98$.}

   \item{Second stabilization mechanism:} what if the bank is out of reserves, but the basecoin price in the liquidity pool is still above the oracle basecoin price? In this case we restore price in the pool by removing liquidity, and there are two possible options here:

   \begin{enumerate}
   \item{Burn:} if the bank if empty, and $r$ is below some threshold, it is allowed to burn \stc{}s in the pool. 
   
   \emph{To instantiate the rule, we propose to burn \stc{} if $r \le 0.95$ for a time period $T_{burn}$. $T_{burn}$ must be quite big, e.g. $1 {\ week}$. We burn enough stablecoins to return to the state of stopped withdrawals, so to $r = 0.98$. After burning, the period is reset, so another burn will be done only after $T_{burn}$ amount of time at least.}

   \item{Extract for the future:} if the bank is empty and $r$ is below some threshold, it is allowed to extract \stc{}s from the pool and lock them in a contract, only to be released in the future when the \stc{} price is above the stablecoin oracle price (i.e. the inverse of the basecoin price peg).
   
   \emph{To instantiate the rule, we propose to extract \stc{}s if $r \le 0.95$ for time period $T_{burn}$ (so the same as in burn). 
   To prioritize extracted funds over arbitrage mint, we do not have delay in releasing contract. We burn enough to return to the state of stopped withdrawals, so close to $r = 0.98$~(exactly, $0.97 \le r \le 0.98$). After extraction, the period is reset, so another extraction will be done after $T_{burn}$ amount of time at least.}
   \end{enumerate}
\end{itemize} 

In addition, we introduce $2\%$ redemption fee for liquidity providers to avoid excessive liquidity hopping.

\section{Stability}
\label{sec:stability}

With second stabilization mechanism~(burning or extraction) in place, the price in the reference market will eventually stabilize. However, for liquidity holders, burning is painful, extraction not as much, but still not desirable. Thus, the \dx{} protocol is trying to avoid it (unlike other protocols, such as Gyroscope or Fei, where redemption rate fails and falls below $1$ immediately when the collateralization ratio falls under $100\%$) by giving markets time to self-stabilize, and then doing interventions. However, this could mean slower stabilization, in comparison with other protocols, but slower stabilization helps the bank to play with possible adversaries in an environment with information assymetry, since humans always have access to more information than the algorithmic bank, and with more flexible decision making as well.

\dx{} is also cautious about minting new \stc{}s, which could mean slow growth of the number of stablecoins issued. This could be inconvenient, especially for big players, but the protocol is focused on stability in the first place. 

Please note that the liquidity pool is disincentivizing massive bank runs due to its constant-product nature. Actually, massive bank runs are simpler for the bank to resolve, in comparison to the slow drain of the worst-case scenario.

\section{Implementations}
\label{sec:implementations}

In this section we consider two concrete implementations of the \dx protocol that were made for the Ergo blockchain, 
a gold-pegged stablecoin named \dxg and a USD-pegged stablecoin named \dxusd. The oracle used in both stablecoins is
implemented using the same underlying oracle protocol.

\subsection{Oracle Pools}
\label{subscec:oracle}

\subsection{\dxg}
\label{subsec:dexygold}

Oracles security is the most important issue for our stablecoin protocol, since this is the only component that does not function autonomously on-chain. Thus, the protocol needs to reward oracles.

Oracle Pools 2.0 framework~\cite{eip23} has support for paying rewards in custom tokens. So we will create GORT (Gold
Oracle Reward Token), which will be used to reward oracles and also developers.

$X$ oracles wil get up to $2 \cdot X$ GORTs per epoch (one datapoint update), per Oracle Pool v2 design~\cite{eip23}~(
$2 \cdot X$ if all the oracles are active). We consider $X = 30$ oracles with one hour long epoch, which means up to $60$
GORTs will be released to oracles.

Any transaction in the protocol that requires the oracle datapoint to be provided as a data-input will charge an oracle fee in \bc{}. This \bc{} will be sent to a BuyBack contract. The BuyBack contract purchases GORT available in a LP on ErgoDex, which is then distributed to the oracle operators.
Thus, the oracle operators receiving GORT tokens can redeem them for the accumulated \bc{} from fees in the same LP.

In addition to oracles, GORT is used to reward developers. To do that, we consider a flat emission contract, which is
releasing the same amount per one hour, so 2 GORT per block during initial period (2 years).

A UI for interacting with the \dxg{} protocol is available at \url{https://crystalpool.cc}.

%TODO: Fix up oracle stuff since it applies to USE + check how DORT

\subsection{\dxusd}
\label{subsec:dexyusd}

%TODO: talk about changes to intervention based on fraction of lp instead of bank reserves + consequence/why.
%TODO: add links to crux finance + mew finance + other UIs???

\subsection{Remarks}
\label{subsec:remarks}
The \dxg{} protocol has been implemented on the Ergo blockchain. The ErgoScript smart-contracts follow a separation-of-concerns design pattern, thus each one handles its own part of the protocol instead of having one giant contract, this results in a total of 18 ErgoScript contracts required for the implementation.
These 18 contracts are orchestrated together with the corresponding off-chain software for the different protocol interactions and trackers necessary to check the status of the oracle and pool peg, along with the relevant delay periods.
The smart-contracts, off-chain code, simulations, and tests can be found in the following repository \url{https://github.com/kushti/dexy-stable}. 

\bibliographystyle{plain}

\bibliography{sources}

\end{document}